\section{Follow-up: where does the leaked signal live?}
\label{sec:followup}

Study~2 established \emph{that} crop-level splitting inflates accuracy and that the inflation is
leakage, but neither registered mechanism test isolated \emph{how} the models use it: the
background route was small and uncertain, and the adversary was evaded for one of its two
architectures. We therefore registered eight further hypotheses (F1--F8) as Amendment~8
(commit \texttt{\FollowupCommit}, \FollowupTime), after the Study~2 analysis and before any follow-up model was trained or any
follow-up analysis was run on existing outputs. They form their own Holm family and do not
alter the Study~2 verdicts. Three arms reuse saved Study~2 artifacts; three train new models
on seeds no earlier arm used.

\paragraph{Slide context in CRIC (F1).} CRIC slide images contain several classes, so a test
cell's training siblings can carry a different label. For every cell, the \emph{context class}
is the majority class of the other cells of its slide. If crop-level models apply a slide-level
label prior, a \emph{discordant} cell (whose class differs from its context class) should
receive more probability on the context class when its slide was in training (P1) than when it
was not (P2). F1 is that within-cell difference, with a cluster bootstrap over slides.

\paragraph{Carrier ablations at test time (F2).} Every 2C model is re-evaluated on its test set
after one of five interventions that keep cell morphology: greyscale, Reinhard stain transfer
to an unseen field of the same class, low-pass filtering, illumination flattening, and
replacement of the background by a constant. F2 asks whether stain transfer shrinks the
leak-minus-clean advantage. A quantitative Grad-CAM analysis~\citep{selvaraju2017gradcam} of the
attribution falling outside the cell mask is reported descriptively.

\paragraph{Carrier removal in training (F3, F4).} The 2C factorial is rerun on ten fresh seeds
under three conditions: the standard recipe, every crop's background masked in training and
testing, and strong optical-density (HED) colour augmentation~\citep{tellez2019quantifying}
during training. F3 asks whether leakage survives with no background at all; F4 whether
colour augmentation shrinks it. The standard condition doubles as an independent replication
of H3.

\paragraph{Three ways of removing the fingerprint (F5--F8).} Each manipulation lowers the
within-class parent recoverability $R$ by a different route, and the shortcut account predicts
that the P1 $-$ P2 gap falls with it. F5 erases, from every saved Phase~1 model, the top-$r$
directions of within-class parent identity with a rank-limited LEACE projection
\citep{belrose2023leace} and refits the linear head ($r\in\{0,4,16,64,256\}$). F6 normalises every
crop's colour statistics toward one reference at increasing strength (Reinhard). F7 repeats the
adversarial arm with a probe-matched adversary (standardised features, one linear layer,
periodic re-initialisation) that cannot hide parent identity from a linear probe while leaving it
readable. F8 pools the dose-response over all three. The H5 regime rule and manipulation check
apply to F5 and F7 and decide which levels enter F8.
